\documentclass[14pt,a4paper]{extarticle}

\usepackage{fontspec}
\usepackage[ukrainian,english]{babel}
\usepackage{amsmath}
\usepackage{amssymb}
\usepackage{geometry}
\usepackage{enumitem}
\usepackage{indentfirst}
\usepackage{microtype}

\geometry{
  left=25mm,
  right=15mm,
  top=20mm,
  bottom=20mm
}

\setmainfont{Times New Roman}
\setlength{\parindent}{1.25cm}
\setlength{\parskip}{0pt}
\setlength{\emergencystretch}{3em}
\linespread{1.0}
\setlist{nosep,leftmargin=1.25cm}

\newcommand{\eng}[1]{\foreignlanguage{english}{#1}}
\newcommand{\sectiontitle}[1]{\par\noindent\textbf{#1}\par}

\begin{document}

\begin{center}
\textbf{Завдання з дисципліни: «Науково-практичні засади управління освітніми проєктами створення програмних продуктів»}
\end{center}

\noindent\textbf{Виконав: аспірант гр. 121-24а, Новосельцев Олександр}

\noindent\textbf{Тема дисертаційного дослідження: «Методи та програмні засоби рельєфного текстурування для систем формування тривимірних графічних зображень».}

\sectiontitle{1. Аналіз теми дисертації з точки зору обґрунтування новизни, гіпотез і впровадження результатів}

Актуальність теми пов'язана з тим, що сучасні інтерактивні графічні системи, ігрові рушії, тренажери, системи віртуальної та доповненої реальності потребують високої деталізації поверхонь за обмежених ресурсів графічного процесора (\eng{GPU}). Збільшення кількості полігонів не завжди є прийнятним, оскільки воно підвищує навантаження на пам'ять, геометричний етап конвеєра та передачу даних. Тому важливими залишаються методи, які створюють ілюзію мікрорельєфу на основі карти висот (\eng{height map}), карти нормалей (\eng{normal map}) і зміщення текстурних координат у шейдері (\eng{shader}).

Дисертаційна тема належить до напряму рельєфного текстурування, зокрема паралакс-мапінгу (\eng{parallax mapping}), пошарового паралакс-мапінгу (\eng{steep parallax mapping}), паралаксного оклюзійного відображення (\eng{parallax occlusion mapping}, POM), рельєфного відображення (\eng{relief mapping}), конусного покрокового відображення (\eng{cone step mapping}) і відображення за відстанями (\eng{distance mapping}). Спільна ідея цих методів полягає в тому, що плоска або низькополігональна поверхня розглядається як базова площина, а видимий рельєф відтворюється через пошук перетину променя з полем висот (\eng{ray-heightfield intersection}) у просторі дотичної (\eng{tangent space}).

Наукова новизна може бути обґрунтована через удосконалення методів і програмних засобів рельєфного текстурування для систем формування тривимірних графічних зображень. Базові методи або мають низьку точність, або потребують багатьох вибірок текстури (\eng{texture samples}), особливо під гострими кутами огляду. Перспективною є гіпотеза про використання попередньо обчислених напрямних значень безпечного кроку (\eng{safe-step}) або значень відстані (\eng{distance}), запакованих у канали RGBA текстури. Це дає змогу частину обчислювальної роботи перенести з етапу виконання шейдера (\eng{runtime}) на етап попередньої обробки (\eng{preprocessing}).

Основна гіпотеза дослідження: якщо в додаткових каналах текстури зберігати інформацію про безпечну відстань або крок для кількох характерних напрямів у просторі дотичної (\eng{tangent space}), то під час рендерингу можна зменшити кількість ітерацій променевого проходження (\eng{ray marching}) порівняно з паралаксним оклюзійним відображенням (POM), зберігаючи прийнятну візуальну якість.

\noindent\textbf{Допоміжні гіпотези:}
\begin{enumerate}
  \item Попередньо обчислені напрямні значення можуть зменшити середню кількість вибірок текстури (\eng{texture samples}) на фрагмент.
  \item Метод буде найбільш корисним для складних карт висот (\eng{height maps}) і для косих кутів огляду, де класичне паралаксне оклюзійне відображення (POM) вимагає більше шарів.
  \item Якість результату залежатиме від кількості напрямів, способу інтерполяції між ними та консервативності значень безпечного кроку (\eng{safe-step}).
  \item Додаткова пам'ять у каналах RGBA може бути виправданою, якщо виграш у продуктивності під час виконання (\eng{runtime performance}) перевищує втрати від попередньої обробки (\eng{preprocessing}) і збільшення обсягу текстурних даних.
\end{enumerate}

Результати дослідження можуть бути впроваджені в освітній процес через лабораторні роботи з комп'ютерної графіки, шейдерного програмування та розробки графічних програмних продуктів. Практичний результат може бути оформлений як проєкт для настільної Unity (\eng{Unity Desktop}) із набором шейдерів (\eng{shaders}), тестових сцен, карт висот, засобів вимірювання частоти кадрів (FPS), часу роботи графічного процесора (\eng{GPU time}) і методичних пояснень.

\sectiontitle{2. Науковці, які працюють в Україні та за кордоном у вибраному напрямі}

В Україні близьким до теми є науковий напрям Вінницького національного технічного університету, пов'язаний із текстуруванням, анізотропною фільтрацією, паралакс-мапінгом (\eng{parallax mapping}) і рельєфним текстуруванням. До релевантних дослідників належать О. Н. Романюк, О. О. Дудник, О. В. Романюк та інші автори публікацій ВНТУ з комп'ютерної графіки. Зокрема, у працях О. Н. Романюка та О. О. Дудника розглядаються особливості паралакс-мапінгу (\eng{parallax mapping}), використання карт відстаней до поверхні (\eng{surface distance maps}), модифіковані MIP-піраміди (\eng{MIP pyramids}) і підвищення продуктивності текстурування.

За кордоном важливими для обраного напряму є такі дослідники: Natalya Tatarchuk, яка відома роботами з паралаксного оклюзійного відображення (\eng{Parallax Occlusion Mapping}) і практичними реалізаціями на графічному процесорі (\eng{GPU}); William Donnelly, який запропонував підхід попіксельного відображення зміщення з функціями відстані (\eng{per-pixel displacement mapping with distance functions}); Jonathan Dummer, автор конусного покрокового відображення (\eng{Cone Step Mapping}); Fabio Policarpo та Manuel M. Oliveira, пов'язані з рельєфним відображенням (\eng{relief mapping}) і розслабленим конусним проходженням (\eng{relaxed cone stepping}); Terry Welsh, який описував паралакс-мапінг з обмеженням зміщення (\eng{offset-limited parallax mapping}); Morgan McGuire, пов'язаний із пошаровим паралакс-мапінгом (\eng{steep parallax mapping}); Yu-Jen Chen і Yung-Yu Chuang, які досліджували анізотропне конусне відображення (\eng{anisotropic cone mapping}).

Ці роботи формують теоретичну і практичну базу для порівняння запропонованого підходу. Українські праці важливі для локального наукового контексту та зв'язку з задачами підвищення продуктивності текстурування, а зарубіжні джерела дають основні алгоритмічні моделі пошуку перетину променя з полем висот (\eng{ray-heightfield intersection}) і прискорення на основі попередньої обробки (\eng{preprocessing-based acceleration}).

\sectiontitle{3. План реалізації гіпотез дослідження}

Перший етап: уточнити домовленість про висоту/глибину (\eng{height/depth convention}), систему координат простору дотичної (\eng{tangent-space coordinate system}) і базову математичну постановку задачі. На цьому етапі потрібно чітко визначити, чи текстура зберігає висоту (\eng{height}), чи глибину (\eng{depth}), як задається напрям погляду (\eng{view direction}), і в якому напрямку зміщуються координати UV.

Другий етап: реалізувати базові методи для порівняння: базовий паралакс-мапінг (\eng{basic parallax mapping}), пошаровий паралакс-мапінг (\eng{steep parallax mapping}), паралаксне оклюзійне відображення (POM), рельєфне відображення (\eng{relief mapping}), конусне покрокове відображення (\eng{cone step mapping}). Усі методи мають працювати на одній тестовій сцені настільної Unity (\eng{Unity Desktop}), з однаковими картами висот (\eng{height maps}), картами нормалей (\eng{normal maps}) і матеріалами.

Третій етап: розробити попередню обробку (\eng{preprocessing}) для отримання додаткових напрямних значень відстані (\eng{distance}) або безпечного кроку (\eng{safe-step}). Для початку доцільно використати невелику двовимірну карту висот (\eng{2D height map}) і вручну перевірити геометричний зміст збережених значень. Потім алгоритм можна автоматизувати для повнорозмірних текстур.

Четвертий етап: реалізувати експериментальний варіант шейдера (\eng{shader variant}), який використовує напрямні дані в RGBA (\eng{RGBA directional data}). Наприклад, канал R може містити базову висоту або глибину (\eng{height/depth}), а G, B, A - значення безпечного кроку (\eng{safe-step}) для трьох обраних косих напрямів у просторі дотичної (\eng{oblique tangent-space directions}).

П'ятий етап: провести порівняльний експеримент за кількістю вибірок текстури (\eng{texture samples}), часом рендерингу, стабільністю на гострих кутах огляду, кількістю артефактів і візуальною подібністю до еталонного методу з великою кількістю кроків.

\sectiontitle{4. Методи дослідження і частина вступу про загальнонаукові та спеціальні методи}

У дисертаційному дослідженні передбачається використання загальнонаукових і спеціальних методів. Методи аналізу та синтезу застосовуються для порівняння існуючих підходів до рельєфного текстурування та формування власного комбінованого підходу. Метод абстрагування використовується для переходу від реальної геометрії поверхні до поля висот (\eng{height field}) у текстурному просторі (\eng{texture space}). Метод моделювання застосовується для побудови математичної моделі взаємодії променя погляду з картою висот. Метод порівняння використовується для оцінювання запропонованого способу відносно паралаксного оклюзійного відображення (POM), рельєфного відображення (\eng{relief mapping}) і конусного покрокового відображення (\eng{cone step mapping}).

До спеціальних методів належать методи комп'ютерної графіки, шейдерного програмування, чисельного пошуку перетину, дискретної обробки карт висот (\eng{height maps}), попередньої обробки текстурних даних (\eng{texture data preprocessing}), профілювання графічного процесора (\eng{GPU profiling}) та експериментального вимірювання продуктивності. Практична реалізація здійснюється в настільній Unity (\eng{Unity Desktop}) із використанням HLSL/ShaderLab. Для оцінювання результатів використовуються метрики кількості вибірок текстури (\eng{texture samples}), часу роботи графічного процесора (\eng{GPU time}), частоти кадрів (FPS), кількості ітерацій, стабільності зміщення UV та візуальної якості.

Фрагмент вступу: Для досягнення мети дослідження використано комплекс загальнонаукових і спеціальних методів. Аналіз наукових джерел дав змогу визначити переваги й обмеження існуючих методів рельєфного текстурування. Математичне моделювання застосовано для опису процесу пошуку перетину видового променя з дискретним полем висот (\eng{height field}) у просторі дотичної (\eng{tangent space}). Експериментальне моделювання в настільній Unity (\eng{Unity Desktop}) використано для перевірки працездатності запропонованого підходу в умовах інтерактивного рендерингу. Методи порівняльного аналізу застосовано для оцінювання продуктивності та якості зображення відносно базових алгоритмів паралакс-мапінгу (\eng{parallax mapping}), паралаксного оклюзійного відображення (POM), рельєфного відображення (\eng{relief mapping}) і конусного покрокового відображення (\eng{cone step mapping}).

\sectiontitle{5. Математичний апарат і математична модель дослідження}

Нехай поверхня полігона має параметризацію через текстурні координати $u, v$. Карта висот задається функцією:
\[
H(u,v), \quad H(u,v) \in [0,1].
\]

Для зручності можна також використовувати глибину:
\[
D(u,v)=1-H(u,v),
\]
де глибина (\eng{depth}) показує, наскільки глибоко промінь зайшов у віртуальний рельєф.

У просторі дотичної (\eng{tangent space}) вісь $z$ відповідає нормалі базової поверхні, а площина $xy$ відповідає напрямам зміни UV. Видовий напрям задається нормалізованим вектором:
\[
V_t=(V_x,V_y,V_z).
\]
Якщо
\[
V_z \approx 0,
\]
камера дивиться майже вздовж поверхні, і UV-зміщення стає великим. Саме тому формули паралакс-мапінгу (\eng{parallax mapping}) містять відношення:
\[
\frac{V_{xy}}{V_z}.
\]

Промінь у текстурно-глибинному просторі (\eng{texture/depth space}) можна описати параметрично:
\[
U(t)=U_0-t\cdot p,
\]
\[
D_{\mathrm{ray}}(t)=t,
\]
де
\[
U_0=(u_0,v_0),
\]
\[
p=h_s\cdot\frac{V_{xy}}{V_z},
\]
$h_s$ - масштаб висоти (\eng{heightScale}), а
\[
t\in[0,1]
\]
є поточною глибиною проходження променя.

Умова перетину для домовленості про глибину (\eng{depth convention}) має вигляд:
\[
D_{\mathrm{ray}}(t)\geq D(U(t)).
\]
Тобто промінь досяг або пройшов глибину, на якій розташована поверхня поля висот (\eng{height field}). Завдання алгоритму полягає в тому, щоб знайти найменше $t$, для якого виконується ця умова. Різні методи відрізняються тим, як саме вони шукають це $t$: рівномірними шарами, грубим пошуком із бінарним уточненням, безпечними кроками на основі конусів (\eng{cone-based safe steps}) або кроками на основі відстаней (\eng{distance-based steps}).

Для експериментального методу вводиться додаткова функція:
\[
S(U,k),
\]
де $k$ - індекс попередньо заданого напряму в просторі дотичної (\eng{tangent space}), а $S$ - значення безпечного кроку (\eng{safe-step}) або відстані (\eng{distance}), збережене в одному з каналів текстури. Тоді крок під час виконання (\eng{runtime step}) може бути не сталим, а залежним від напряму огляду:
\[
\operatorname{step}(U,V_t)=\operatorname{interpolate}(S(U,k),V_t).
\]
Мета полягає в тому, щоб крок був достатньо великим для прискорення пошуку, але не призводив до пропуску першого перетину.

\sectiontitle{6. Конкретизація математичної моделі для організації експерименту}

Для експерименту обирається однакова базова модель: площина або куб у Unity, на яку накладаються текстура базового кольору (\eng{albedo texture}), карта нормалей (\eng{normal map}) і карта висот/глибин (\eng{height/depth map}). Усі розрахунки перетину виконуються в просторі дотичної (\eng{tangent space}). Це дає змогу порівнювати методи незалежно від положення об'єкта у світовому просторі (\eng{world space}).

Фіксуються такі вхідні параметри: роздільна здатність (\eng{resolution}) карти висот, масштаб висоти (\eng{heightScale}), кількість шарів POM, кількість кроків уточнення (\eng{refinement steps}) для рельєфного відображення (\eng{relief mapping}), кількість напрямів для напрямної відстані (\eng{directional distance}), формат RGBA текстури, кут огляду камери, освітлення, роздільна здатність екрана.

Еталонним результатом може бути паралаксне оклюзійне відображення (POM) або рельєфне відображення (\eng{relief mapping}) із великою кількістю кроків, наприклад 64 або 128 вибірок (\eng{samples}). Експериментальний метод порівнюється з еталоном за різних кутів огляду: 0--20 градусів, 20--45 градусів, 45--70 градусів і 70--85 градусів відносно нормалі поверхні.

Вимірювані показники: середня кількість вибірок (\eng{samples}) на фрагмент, час роботи графічного процесора (\eng{GPU time}), частота кадрів (FPS), наявність артефактів перескоку (\eng{overshoot artifacts}), ковзання/розтягування UV (\eng{sliding/stretching UV}), стабільність на межах рельєфу, візуальна різниця з еталоном. Для якісного аналізу можна зберігати знімки екрана (\eng{screenshots}) із однакових позицій камери (\eng{camera positions}).

\sectiontitle{7. План експерименту}

\begin{enumerate}
  \item Підготувати сцену Unity (\eng{Unity scene}) з площиною, кубом і набором матеріалів.
  \item Підготувати кілька карт висот (\eng{height maps}): проста синусоїдна, цегляна/кам'яна, різка ступінчаста, шумова.
  \item Реалізувати режим налагоджувальної візуалізації шейдера (\eng{shader debug mode}) для напряму погляду в просторі дотичної (\eng{tangent-space view direction}).
  \item Реалізувати базовий паралакс-мапінг (\eng{basic parallax mapping}) і перевірити домовленість про знак (\eng{sign convention}).
  \item Реалізувати POM із фіксованою та залежною від кута огляду (\eng{view-angle-dependent}) кількістю шарів.
  \item Реалізувати рельєфне відображення (\eng{relief mapping}): грубий пошук (\eng{coarse search}) плюс бінарне уточнення (\eng{binary refinement}).
  \item Реалізувати конусне покрокове відображення (\eng{cone step mapping}) або спрощений базовий метод на основі конусів (\eng{cone-based baseline}).
  \item Реалізувати попередню обробку (\eng{preprocessing}) для напрямної RGBA карти відстаней/безпечних кроків (\eng{directional RGBA distance/safe-step map}).
  \item Реалізувати експериментальний шейдер (\eng{shader}), який використовує додаткові канали RGBA.
  \item Провести серію вимірювань для різних кутів огляду, масштабу висоти (\eng{heightScale}) і типів карт висот (\eng{height maps}).
  \item Зібрати таблицю результатів: метод (\eng{method}), середня кількість вибірок (\eng{average samples}), час роботи графічного процесора (\eng{GPU time}), частота кадрів (FPS), артефакти (\eng{artifacts}), вартість пам'яті (\eng{memory cost}).
  \item Порівняти результати з еталоном і сформулювати висновки щодо підтвердження або спростування гіпотез.
\end{enumerate}

\sectiontitle{8. Очікувані результати експерименту}

Очікується, що базовий паралакс-мапінг (\eng{basic parallax mapping}) буде найшвидшим, але матиме помітні артефакти ковзання і неправильного перетину, особливо під гострими кутами. Паралаксне оклюзійне відображення (POM) має дати кращу візуальну якість, але потребуватиме більше вибірок текстури (\eng{texture samples}). Рельєфне відображення (\eng{relief mapping}) може дати точніше уточнення перетину, але його якість залежатиме від того, чи знайдено правильний інтервал для бінарного пошуку (\eng{binary search}). Конусне покрокове відображення (\eng{cone step mapping}) має зменшити кількість кроків під час виконання (\eng{runtime steps}) завдяки попередній обробці (\eng{preprocessing}), але потребуватиме зберігання додаткових даних.

Для експериментального підходу з напрямними відстанями в RGBA (\eng{RGBA directional distance}) очікується потенційне зменшення середньої кількості кроків порівняно з POM. Найбільший ефект може проявитися на картах висот (\eng{height maps}) із вираженим рельєфом і при косих кутах огляду. Водночас можливі ризики: перескок через перший перетин (\eng{overshoot}), некоректна інтерполяція між напрямами, збільшення пам'яті, складність попередньої обробки (\eng{preprocessing}) і поява артефактів на різких межах.

Якщо експеримент покаже, що метод зменшує кількість вибірок (\eng{samples}) без суттєвої втрати якості, це підтвердить доцільність подальшого розвитку гіпотези. Якщо артефакти будуть значними, результат усе одно буде корисним, бо дасть змогу визначити обмеження представлення напрямного безпечного кроку (\eng{directional safe-step representation}) і уточнити математичну модель.

\sectiontitle{9. План освітнього проєкту за темою дисертації}

Назва освітнього проєкту: «Рельєфне текстурування в Unity: від базового паралакс-мапінгу (\eng{basic parallax mapping}) до прискорених методів пошуку перетину променя з полем висот (\eng{ray-heightfield intersection})».

Мета проєкту: створити навчальний модуль для студентів комп'ютерної графіки, який пояснює геометричний зміст рельєфного текстурування та дає практичні навички реалізації шейдерів у Unity.

Цільова аудиторія: студенти спеціальностей комп'ютерні науки, програмна інженерія, комп'ютерна інженерія, а також розробники-початківці, які вивчають рендеринг у реальному часі (\eng{real-time rendering}).

Очікувані освітні результати: студент пояснює різницю між відображенням нормалей (\eng{normal mapping}), паралакс-мапінгом (\eng{parallax mapping}), паралаксним оклюзійним відображенням (POM) і рельєфним відображенням (\eng{relief mapping}); розуміє простір дотичної (\eng{tangent space}) і базис дотична-бінормаль-нормаль (TBN); реалізує простий шейдер HLSL (\eng{HLSL shader}); проводить вимірювання продуктивності; аналізує компроміс (\eng{trade-off}) між якістю, кількістю вибірок (\eng{samples}), пам'яттю і попередньою обробкою (\eng{preprocessing}).

\noindent\textbf{Структура проєкту:}
\begin{enumerate}
  \item Лекція 1: поля висот (\eng{height fields}), простір дотичної (\eng{tangent space}) і базова ідея паралакс-мапінгу (\eng{parallax mapping}).
  \item Лабораторна робота 1: візуалізація напряму погляду в просторі дотичної (\eng{tangent-space view direction}) у Unity.
  \item Лекція 2: пошаровий паралакс-мапінг (\eng{steep parallax mapping}) і POM як променеве проходження (\eng{ray marching}) у текстурному просторі (\eng{texture space}).
  \item Лабораторна робота 2: реалізація POM і налагоджувальна візуалізація шляху UV (\eng{UV path}).
  \item Лекція 3: рельєфне відображення (\eng{relief mapping}), бінарне уточнення (\eng{binary refinement}) і конусне покрокове відображення (\eng{cone step mapping}).
  \item Лабораторна робота 3: порівняння POM і рельєфного відображення (\eng{relief mapping}).
  \item Проєктна робота: власний варіант шейдера (\eng{shader variant}) або попередня обробка текстурних даних (\eng{preprocessing texture data}).
  \item Підсумкова робота: звіт із результатами вимірювань і висновками.
\end{enumerate}

Програмний продукт освітнього проєкту: пакет Unity (\eng{Unity package}) або репозиторій із тестовою сценою, варіантами шейдерів (\eng{shader variants}), картами висот (\eng{height maps}), інструкціями до лабораторних робіт і шаблоном звіту. Такий проєкт може бути використаний у курсах комп'ютерної графіки, розробки ігрових рушіїв, програмування графічного процесора (\eng{GPU programming}) та освітніх проєктів створення програмних продуктів.

\sectiontitle{Висновок}

Запропонована тема дисертації має наукову і практичну значущість, оскільки спрямована на розвиток методів та програмних засобів рельєфного текстурування для систем формування тривимірних графічних зображень. Її новизна може полягати в поєднанні класичних методів пошуку перетину променя з полем висот (\eng{ray-heightfield intersection}) із додатковими попередньо обчисленими напрямними даними в текстурних каналах. Практична реалізація в настільній Unity (\eng{Unity Desktop}) дає змогу перевірити гіпотези експериментально, а результати можуть бути безпосередньо використані в освітньому процесі як навчальний модуль із комп'ютерної графіки та шейдерного програмування.

\sectiontitle{Використані та орієнтовні джерела}

\begin{enumerate}
  \item О. Н. Романюк, О. О. Дудник. Використання модифікованої MIP-піраміди для підвищення продуктивності паралакс-мапінгу (\eng{Parallax Mapping}). Вісник Вінницького політехнічного інституту.
  \item О. О. Дудник. Дисертація: методи та засоби підвищення реалістичності і продуктивності текстурування в системах комп'ютерної графіки. ВНТУ, 2018.
  \item N. Tatarchuk. Практичне динамічне паралаксне оклюзійне відображення (\eng{Practical Dynamic Parallax Occlusion Mapping}). SIGGRAPH Sketches, 2005.
  \item N. Tatarchuk. Динамічне паралаксне оклюзійне відображення з наближеними м'якими тінями (\eng{Dynamic Parallax Occlusion Mapping with Approximate Soft Shadows}). SI3D, 2006.
  \item W. Donnelly. Попіксельне відображення зміщення з функціями відстані (\eng{Per-pixel displacement mapping with distance functions}). GPU Gems 2, 2005.
  \item J. Dummer. Конусне покрокове відображення: ітеративний алгоритм пошуку перетину променя з полем висот (\eng{Cone Step Mapping: An Iterative Ray-Heightfield Intersection Algorithm}), 2006.
  \item F. Policarpo, M. M. Oliveira. Розслаблене конусне проходження для рельєфного відображення (\eng{Relaxed Cone Stepping for Relief Mapping}). GPU Gems 3, 2007.
  \item M. McGuire. Пошаровий паралакс-мапінг (\eng{Steep Parallax Mapping}), 2005.
  \item T. Welsh. Паралакс-мапінг з обмеженням зміщення (\eng{Parallax Mapping with Offset Limiting}), 2004.
  \item Yu-Jen Chen, Yung-Yu Chuang. Анізотропне конусне відображення (\eng{Anisotropic Cone Mapping}), 2009.
\end{enumerate}

\end{document}
